\chapter{Homecoming}
Time slipped like sand through his fingers as Shiva entered ninth standard. This was the age when the school van was left behind, when the small magic of being driven gave way to something earned—a bicycle of his own. He got a geared cycle, its brakes betraying him on the very first day.

The thirty-minute ride from school to home felt like a declaration of independence, even as the Indian summer’s loo lashed against his face. Pedaling along the same road, beside the vast, open riverbed, he overtook the familiar school van. Watching it move past him—slow, enclosed, unchanged—he felt something settle inside him. He was no longer a boy being carried home. He was moving forward on his own.

Shreya’s call had woven itself into the fabric of his afternoons, a ritual as certain as the setting sun. Returning home no longer felt like an ending, but a prelude to the moment he could finally unload the day’s weight into the hum of the telephone. Most evenings were smooth, a quiet exchange of ordinary moments, but sometimes, the air between them grew thin with a sudden, unexpected tension. One such evening, Shiva sat on the terrace, his eyes tracing the familiar flocks of egrets as they painted the deepening sky with streaks of white, the phone pressed tight to his ear. 

``We’re so perfect together, you know?'' Shreya’s voice drifted through the receiver, buoyant and thick with an emotion she didn't try to hide. ``I showed our picture to my friends at coaching today. They all said we’re... well, they said we’re made for each other.'' 

Shiva felt a slow smile tug at his lips. ``I think they’re right,'' he murmured. ``We’re like Radha and Krishna. Destined.'' 

There was a sharp silence on the other end, as if the connection had suddenly frayed. ``No,'' Shreya said, her voice dropping into a register that made Shiva sit up. ``I’m not your Radha.'' 

The world seemed to tilt slightly on its axis. ``What is it?'' Shiva asked, his heart hammering against his ribs. ``Did I say something wrong?'' 

``I don’t want to be Radha,'' she said firmly, though the softness returned to her tone. ``I want to be Rukmini. I want to be the one who stays, the one who is truly his.'' 

A quiet understanding washed over Shiva, as warm as the fading sunlight. ``I see,'' he said, the tension leaving his shoulders. ``Names are just words, Shreya. You’re everything to me. No single story could ever hold what I feel for you.'' 

``I know,'' she whispered, and he could almost hear the smile returning to her face. ``I feel the same.'' 

As the new academic session gathered momentum, a new rhythm took over the village. One by one, Shiva’s friends began disappearing into the local coaching centers, their cycles parked in crowded rows outside small, rented rooms. Shiva didn’t want to be the one left behind, anchored to the porch while the world moved on. He had his geared bicycle now—his own wings. He pleaded with Nani for days, promising to be careful, to watch the highway, to never race. She hesitated, her worries etched into the slow way she turned her prayer beads, but eventually, his persistence won. She gave her blessing, though it came with a long list of warnings that he barely heard in his excitement. 

The coaching shifted the very foundations of his mornings. He was out of bed by six, the air still cool and smelling of damp earth. While the rest of the house was still waking, Mami would already be in the kitchen, the hiss of the stove a steady companion as she prepared a breakfast that was always too large for the hour. She moved with a quiet efficiency, ensuring he left with a full stomach, while Nani tucked packets of biscuits into his bag—remnants of a care that refused to believe he didn't need a midday meal. By 7:30, he was on the road, pedaling toward the 8:00 AM class. It was a tight, perfect loop: an hour of coaching, then a quick dash across the town to reach the school gates by 9:30. 

In the cramped coaching room, Shiva often found himself listening to explanations for things he had already mastered in the quiet of his own room. His marks in the weekly tests were consistently at the top, prompting the teacher to occasionally look up from his register and ask why he even bothered coming. Shiva would only offer a shy, practiced smile. ``I just like the revision, sir,'' he’d say. In truth, it was more than that—it was the proximity to the life his friends were living. Vijay sat beside him, a familiar anchor in the sea of new faces, their shared silence a comfort during the long lectures. 

The highway, however, held shadows that the village lanes did not. One morning, the usual blur of the roadside was broken by a cluster of people gathered near the opposite shoulder. A car lay there, its frame twisted into an unnatural shape, glass glinting like scattered diamonds on the asphalt. It was a sight common enough on that stretch of road to be greeted with a grim, passing glance. Shiva pedaled past, his mind already on the day’s lessons, unaware of how close the shadow had fallen. It was only later, sitting on the narrow wooden bench in the coaching center, that the reality settled in. 

``A man from my village,'' Vijay whispered, his face unusually pale. He had just heard the news from his uncle. ``He was hit this morning. Right there on the highway.'' 

A cold knot tightened in Shiva’s stomach, a sharp, icy contrast to the humid air outside. The image of the mangled car returned to him—the twisted metal, the indifferent crowd, the glint of shattered glass on the hot asphalt. ``Was he... was he hurt badly?'' he asked, his voice sounding thin even to his own ears. 

``I don't know much,'' Vijay replied, his words barely audible over the hum of the classroom. ``I was leaving for the day when the news reached us. Everything was a blur.''

The silence that followed was heavy, a thick layer of unease that the ceiling fan’s rhythmic drone couldn’t cut through. Shiva looked down at his hands, resting on the scarred wooden desk—the same hands that, only an hour ago, had gripped his bicycle handlebars with the confidence of someone who owned the road.

The lecture began shortly after, the teacher’s voice a distant murmur as the class dissolved into a blink of notes and clock-watching. When the bell finally rang, the Tridev split apart—Shiva and Vijay heading toward the school on their cycles, while they called out to Rajveer on the phone, their voices lost in the rush of traffic. 

The days settled back into their familiar rhythm, a cycle of morning air and afternoon heat. And though Shiva had not spoken a word to Gauri, nor let her name slip into his conversations with anyone else, his eyes never failed to find her. She was a silent anchor in his classroom, a presence he tracked without needing to look. 

But lately, her seat had been empty. One day of absence turned into two, then three, until the classroom felt strangely hollow, as if a vital piece of the day-to-day architecture had been removed. Finally, the silence grew too loud to ignore. Shiva waited until a quiet moment in the corridor before approaching Divya. 

``Where is Gauri these days?'' he asked, trying to keep his voice as casual as the dust on the floor. ``She hasn't been in class for a while.'' 

Divya paused, a slow, knowing smirk spreading across her face. ``Oho, so you do notice her when you think we’re not looking...'' 

``It’s not like that,'' Shiva stammered, the heat rising in his cheeks despite the cool corridor breeze. ``I just... I noticed the empty bench, that's all. It’s hard not to notice.'' 

``You never asked about me when I was away, though,'' Divya teased, her eyes glinting with a familiar, sharp mischief. 

``I do! I would!'' Shiva said, his defense stumbling over itself. ``You’re such a nuisance, Divya. Why can’t you just give me a straight answer?'' 

Divya dispensed with the teasing and laughed, the sound bright and effortless. ``Relax, Shiva. She’s just a little unwell. A change in the weather, probably. She’ll be back as soon as the fever breaks.'' 

The words were simple, yet they acted like a release valve, letting the pressure of the last few days escape. Shiva nodded, a quiet sense of relief settling over him like the early evening light.

